\documentclass[11pt]{article}
\usepackage[a4paper,margin=25mm]{geometry}
\usepackage{amsmath,amssymb,amsthm}
\usepackage[hidelinks]{hyperref}
\emergencystretch=2em
\newtheorem{theorem}{Theorem}
\newtheorem{lemma}{Lemma}
\newcommand{\Z}{\mathbb Z}
\newcommand{\C}{\mathbb C}
\newcommand{\supp}{\operatorname{supp}}
\title{A support-product minimizer whose supports are not difference sets\\
\large Disproof of Conjecture 00000009278}
\author{gaochengzhecpu}\date{}
\begin{document}\maketitle
\begin{abstract}
For complex functions on $G=\Z/4\Z$, the smallest nonzero Fourier
support product is $4$. The indicator of $H=\{0,2\}$ attains this
minimum, and both its support and its Fourier support equal $H$.
This set is not a difference set. Thus the asserted difference-set
classification of extremals is false, although the usual uncertainty
lower bound remains valid. We prove global minimality over all
nonzero complex functions, not just optimality within a test family.
\end{abstract}

\section*{Conventions and the claim being refuted}
Identify the dual of $G=\Z/4\Z$ with $G$ by its standard characters.
For $f:G\to\C$, use the unnormalized Fourier transform
\[
 \widehat f(k)=\sum_{x\in G}f(x)e^{-2\pi i xk/4},\qquad
 \supp f=\{x\in G:f(x)\ne0\}.
\]
Changing to the normalized transform does not change either support.
The support product is $|\supp f|\,|\supp\widehat f|$, for $f\ne0$.
An extremal here means a function attaining its global minimum.

A subset $D\subseteq G$ is a difference set if some integer
$\lambda\ge0$ satisfies
\[
 \#\{(a,b)\in D^2:a-b=g\}=\lambda
       \quad\text{for every }g\in G\setminus\{0\}.
\]
This definition even allows degenerate parameter choices. We disprove
the source's assertion that support-product extremals are difference
sets: our minimizing function has neither a difference-set support
nor a difference-set Fourier support. A cyclic finite abelian group
already lies within the stated finite-group scope.

\section*{The actual minimum is four}
The elementary inversion identity is
\[
 \widehat{\widehat f}(x)=4f(-x).
\]
Indeed, expand the two sums: the inner sum over $k$ is $4$ if the
frequency is zero modulo $4$, and $0$ otherwise. In particular
$f\ne0$ implies $\widehat f\ne0$.

\begin{lemma}
Every nonzero function $f:G\to\C$ satisfies
$|\supp f|\,|\supp\widehat f|\ge4$.
\end{lemma}
\begin{proof}
Both support sizes are positive integers. If each is at least $2$,
their product is at least $4$.
If $\supp f=\{j\}$, then
\[
 \widehat f(k)=f(j)e^{-2\pi i jk/4}\ne0
           \quad\text{for every }k\in G,
\]
so the sizes are $1$ and $4$.
If $\supp\widehat f$ is a singleton, the same argument applied to
$\widehat f$ shows that $\widehat{\widehat f}$ is nonzero everywhere.
Inversion then shows that $f$ is nonzero everywhere. The sizes are
$4$ and $1$. These cases exhaust all possibilities.
\end{proof}

\section*{An extremal with two nondifference-set supports}
Take the proper nontrivial subgroup $H=\{0,2\}$ and $f=1_H$.
Directly from the transform definition,
\[
 \widehat f(k)=1+e^{-\pi i k}=1+(-1)^k
       =\begin{cases}2,&k\in\{0,2\},\\0,&k\in\{1,3\}.
         \end{cases}
\]
Thus $f=(1,0,1,0)$, $\widehat f=(2,0,2,0)$, and
\[
 \supp f=\supp\widehat f=H,\qquad
 |\supp f|\,|\supp\widehat f|=2\cdot2=4.
\]
The lemma proves that this is an actual global minimum.

However, the differences of the four ordered pairs in $H^2$ are
\[
 0-0=0,\quad 0-2=2,\quad 2-0=2,\quad 2-2=0
                   \quad(\bmod 4).
\]
Hence nonzero difference $1$ occurs $0$ times and nonzero difference
$2$ occurs $2$ times. No common $\lambda$ exists.

\begin{theorem}
There is a nonzero complex function on a finite group attaining the
minimum Fourier support product such that neither support is a
difference set. Consequently the extremal-classification assertion
in Conjecture 00000009278 is false.
\end{theorem}
\begin{proof}
The function $1_H$ constructed above has all the stated properties.
\end{proof}

\section*{Formal correspondence and reproduction}
The Lean project works with the actual group \texttt{ZMod 4}, complex
values, and Mathlib's standard \texttt{ZMod.dft} linear equivalence.
It uses the library's Fourier inversion theorem, rather than replacing
the transform by a custom table. The support predicate is precisely
nonzero function values. The theorem \texttt{support\_product\_bound}
quantifies over every nonzero complex function on the group.

The project constructs $H$ as both a finite set and an additive
subgroup. The theorem \texttt{fourier\_witness} evaluates the actual
transform and proves $\widehat{1_H}=2\,1_H$. Character values are
derived from the standard character's homomorphism and injectivity
properties. Both support identifications and the product $4$ are
proved. The difference-set predicate counts actual ordered pairs
with the specified group difference; the counts $0$ and $2$ are
kernel checked. The final theorem \texttt{counterexample} asserts
nonzeroness, global minimality, and the failure of the difference-set
property for both supports.

Lean 4.19.0 and Mathlib revision\\
\texttt{c44e0c8ee63ca166450922a373c7409c5d26b00b}
are pinned. In \texttt{lean/}, run
\begin{verbatim}
lake exe cache get
lake build
lake env lean -DwarningAsError=true Main.lean
\end{verbatim}
The cache command obtains the pinned official dependency artifacts
if absent. The submitted project is rebuilt in a fresh directory.
No auxiliary numerical program, extra axiom, proof placeholder, or
external computational oracle is used. Run \texttt{tectonic main.tex}
to regenerate the PDF. The author performed a separate adversarial
self-review; no independent reviewer or subagent was used.

\section*{Source}
The Last Math Competition, Conjecture 00000009278.
The exact bilingual statement is preserved in \texttt{SOURCE.md};
its upstream revision, source hash, actual build logs, axiom
dependencies and review records accompany this submission.
\end{document}
